\newcommand{\decisionproblem}[2]
{
	\begin{itemize}
		\item {\bf Input.} #1
		\item {\bf Output.} #2
	\end{itemize}
}

\newcommand{\freevar}[1]{\mathrm{var}(#1)}

\newcommand{\fdp}{\textsc{fdp}\xspace}

\newcommand{\pspace}{{\sc{PSpace}\xspace}}
\newcommand{\indvec}[1]{\atoms^{\langle #1 \rangle}}
 \newcommand{\dimension}[1]{\mathrm{dim}(#1)}
\newcommand{\leastsup}[1]{\mathrm{sup}(#1)}

\newcommand{\syntdim}[1]{\mathrm{dim}(#1)}
\newcommand{\syntsize}[1]{|#1|}
\newcommand{\semdim}[1]{\mathrm{\red{dim}}(#1)}
\newcommand{\semsize}[1]{\red|#1\red|}

\newcommand{\orbitsize}[2]{|#1|_{#2}}
%\newcommand{\saut}[1]{\mathrm{G}_{#1}}
%\newcommand{\aut}{\mathrm{G}}
%\newcommand{\secfm}{\bfseries{\scshape{fm\ }}}
%\newcommand{\fm}{\textsc{fm}\ }
%\newcommand{\tneq}{\(\neq\)}
\newcommand{\tleq}{\(\leq\)}
% \newcommand{\tin}{\(\in\)}
% \newcommand{\tnotin}{\(\notin\)}
% \newcommand{\tcup}{\(\cup\)}
% \newcommand{\temptyset}{\(\emptyset\)}
% \newcommand{\tcap}{\(\cap\)}
% \newcommand{\ttimes}{\(\times\)}

\newcommand{\approxi}[2]{#1^{(#2)}}

\newcommand{\cellc}[2]{\mathsf{cell}_{#1}(#2)}
\newcommand{\tatoms}{\(\ \!\!\!\!\atoms\)}

% \newcommand{\booltrue}{\ensuremath{\mathtt{true}}}
% \newcommand{\boolfalse}{\ensuremath{\mathtt{false}}}
\newcommand{\tset}[1]{\{#1\}}

\newcommand{\settype}{\mathtt{set}}
\newcommand{\atomtype}{\mathtt{atom}}

\newcommand{\red}[1]{{\color{red}#1}}
\newcommand{\blue}[1]{{\color{blue}#1}}

\newcommand{\exercisepart}{

\bigskip

{\noindent \large \bf Exercises}


}
\newcommand{\setexpr}[3]{\set{#1\ |\ \text{ for }#2 \in \atoms \text{ such that }#3}}
\newcommand{\setexprtup}[4]{\set{#1\ |\ \text{ for }#2 \in \atoms^{#3} \text{ such that }#4}}
\newcommand{\setexprtrue}[2]{\set{#1\ |\ \text{ for }#2 \in \atoms}}
\newcommand{\setexprtuptrue}[3]{\set{#1\ |\ \text{ for }#2 \in \atoms^{#3}}}

\newcommand{\symbolpushbis}[2]{\symbolpush{#1 \ : \ #2}}
\newcommand{\symbolpush}[1]{\underline{#1}}

\newcommand{\aequiv}[1]{\stackrel{#1}\sim}
\newcommand{\locations}{\mathsf{Loc}}

\newcommand{\bibnotes}[1]{
\bigskip
{\footnotesize
\noindent {\bf Bibliographic notes for Section~\arabic{section}. }
#1}
}

% these two macros are for delayed exercises, I don't use them now
% \renewcommand{\ExerciseHeader}{\noindent \textbf{\ExerciseName~\ExerciseHeaderNB\ExerciseHeaderTitle
% \ExerciseHeaderOrigin. }}
% \renewcommand{\AnswerHeader}{\noindent {\textbf{Solution to \ExerciseName\ \ExerciseHeaderNB.\\}}}

\newcounter{exercisecounter}
\newenvironment{mexercise}{
\refstepcounter{exercisecounter}

\smallskip\noindent{\textbf{{Exercise \arabic{exercisecounter}. }}}}{
}






\newcommand{\pv}[1]{{\mathtt{#1}}}

% the original picture files, from illustrator
\newcommand{\mypic}[1]{
\begin{center}
	\includegraphics[page=#1,scale=0.5]{atom-pics}
\end{center}
}

% for pictures with figma
\newcommand{\mypicb}[1]{
\begin{center}
	\includegraphics[page=#1,scale=0.22]{atom-book}
\end{center}
}

%\newcommand{\notin}{\not \in}

\newcommand{\sets}{\mathsf{set}}
\newcommand{\setbuil}{\mathsf{setbuil}}
\newcommand{\defset}{\mathsf{hdef}}
\newcommand{\hofset}{\mathsf{hof}}

\newcommand{\setbuild}[2]{\set{\ #1 \ | 
\begin{tabular}{l}
	#2
\end{tabular}}}

% the tabular in the above macro can make it too wide, hence this one
\newcommand{\setbuildoneline}[2]{\set{\ #1 \ | 
\text{ #2}\ }}

\newcommand{\myunderbrace}[2]{\underbrace{#1}_{\mathclap{\text{\scriptsize 
\begin{tabular}{c}
	#2
\end{tabular} }}}}
\newcommand{\myoverbrace}[2]{\overbrace{#1}^{\mathclap{\text{\scriptsize 
\begin{tabular}{c}
	#2
\end{tabular} }}}}

\newcommand{\Bool}{2}
\newcommand{\eve}{\text{Eve}}
\newcommand{\adam}{\text{Adam}}
\newcommand{\john}{\text{John}}
\newcommand{\tom}{\text{Tom}}
\newcommand{\jim}{\text{Jim}}
\newcommand{\mary}{\text{Mary}}

\newcommand{\deflift}{{\mathsf{def}}}
\newcommand{\tti}{{\mathtt I}}
\newcommand{\ttj}{{\mathtt J}}
\newcommand{\qatom}{(\mathbb Q, <)}
\newcommand{\sem}[1]{[\![#1]\!]}
\newcommand{\fraisse}{Fra\"iss\'e } 

\newcommand{\saut}[1]{\mathrm{G}_{#1}}
\newcommand{\aut}{\mathrm{G}}
\newcommand{\secfm}{\bfseries{\scshape{fm\ }}}
%\newcommand{\fm}{\textsc{fm}\ }
\newcommand{\tneq}{\(\neq\)}
\newcommand{\tin}{\(\in\)}
\newcommand{\tnotin}{\(\notin\)}
\newcommand{\tcup}{\(\cup\)}
\newcommand{\temptyset}{\(\emptyset\)}
\newcommand{\tcap}{\(\cap\)}
\newcommand{\ttimes}{\(\times\)}

\newcommand{\structclass}{\mathscr A}
\newcommand{\structclassb}{\mathscr B}
\newcommand{\structclassh}{\mathscr H}
\newcommand{\structa}{\mathbb A}
\newcommand{\structb}{\mathbb B}
\newcommand{\structc}{\mathbb C}
\newcommand{\structh}{\mathbb H}

\newcommand{\atoms}{\mathbb A}
\newcommand{\field}{\mathbb F}
\newcommand{\universe}[1]{\mathrm{universe(#1)}}
\newcommand{\Fff}{\mathscr F}
\newcommand{\limmodel}{\underline \structa}
\newcommand{\arity}{\mathrm{arity}}
\newcommand{\lincomb}{\mathrm{Lin}}
\newcommand{\length}{\mathrm{len}}

% various kinds of function spaces
\newcommand{\lineqfun}[2]{ #1 \underset{\text{lineq}}{\longrightarrow} #2}
\newcommand{\linfsfun}[2]{ #1 \underset{\text{linfs}}{\longrightarrow} #2}
\newcommand{\eqfun}[2]{ #1 \underset{\text{eq}}{\longrightarrow} #2}
\newcommand{\fsfun}[2]{ #1 \underset{\text{fs}}{\longrightarrow} #2}

\newcommand{\generate}[1]{\langle #1 \rangle}

\newcommand{\atoma}{{a}}
\newcommand{\atomb}{{b}}
\newcommand{\atomc}{{c}}
\newcommand{\atomd}{{d}}

\newcommand{\atomq}{{q}}
\newcommand{\atomzero}{{\underline 0}}
\newcommand{\atomone}{{\underline 1}}
\newcommand{\atomtwo}{{\underline 2}}
\newcommand{\atomthree}{{\underline 3}}

 
\newcommand{\atomseta}{{A}}
\newcommand{\atomsetb}{{B}}
\newcommand{\atomsetc}{{C}}
\newcommand{\atomsets}{{S}}
\newcommand{\atomsett}{{T}}

\newcommand{\Qfin}{Q_{\mathrm{fin}}}
\newcommand{\Afin}{A_{\mathrm{fin}}}

\newcommand{\bind}{\nabla}
\newcommand{\powerset}{{\mathsf P}}
\newcommand{\pfin}{{\mathsf P}_{\text{fin}}}
\newcommand{\decproblem}[3]
{\begin{center}
\fbox{
\begin{tabular}{rl}
\textsc{Problem}: & #1 \\
\textsc{Input}: & #2\\
\textsc{Output}: & #3
\end{tabular}
}
\end{center}}

\newcounter{openquestioncounter}
\newenvironment{openquestion}{
\medskip

\refstepcounter{openquestioncounter}
\smallskip\noindent{\textbf{{Open question \arabic{openquestioncounter}. }}}}{\hfill $\Box$
\medskip 
}







\newcommand{\eqdef}{\stackrel {\mathrm{def}} = }

\newcommand{\Ii}{\mathcal I}